% Artículo IV: la conjetura CST de Terras por número de rachas (Ocelotas).
% Versión del 12-10-2026 (81 rachas con pesos; 42 con censo propio hasta 3e11).
% El borrador del 04-10 (24 y 77 rachas) se envió ese día a Rozier-Terracol y a Hercher.
% Pendiente: lectura humana; MathSciNet no consultado (zbMATH Open sí, 12-10).
% Repositorio citado en "Data availability" en el commit ccedd6cfd3ff (push del 12-10).
% Compilar con: tectonic -X compile articulo_IV_ocelotas.tex   (o pdflatex, dos pasadas)
\documentclass[11pt]{amsart}
\usepackage{iftex}
\ifPDFTeX
  \usepackage[utf8]{inputenc}
  \usepackage[T1]{fontenc}
\fi
\usepackage{amssymb}
\usepackage{booktabs}
\usepackage{array}
\usepackage{fancyvrb}
\usepackage{newunicodechar}
\usepackage[hidelinks]{hyperref}
\hypersetup{pdftitle={On Terras's coefficient stopping time conjecture for Collatz sequences with few local maxima},pdfauthor={Omar Javier Said Duran}}
\usepackage{xurl}
\newunicodechar{ℕ}{\ensuremath{\mathbb{N}}}
\newunicodechar{ℝ}{\ensuremath{\mathbb{R}}}
\newunicodechar{→}{\ensuremath{\to}}
\newunicodechar{∀}{\ensuremath{\forall}}
\newunicodechar{∃}{\ensuremath{\exists}}
\newunicodechar{∧}{\ensuremath{\wedge}}
\newunicodechar{∨}{\ensuremath{\vee}}
\newunicodechar{≤}{\ensuremath{\le}}
\newunicodechar{≠}{\ensuremath{\ne}}
\newunicodechar{…}{\ensuremath{\dots}}

\newtheorem{mainthm}{Theorem}
\renewcommand{\themainthm}{\Alph{mainthm}}
\newtheorem{maincor}[mainthm]{Corollary}
\newtheorem{theorem}{Theorem}[section]
\newtheorem{lemma}[theorem]{Lemma}
\newtheorem{proposition}[theorem]{Proposition}
\newtheorem{corollary}[theorem]{Corollary}
\theoremstyle{definition}
\newtheorem{definition}[theorem]{Definition}
\newtheorem{example}[theorem]{Example}
\theoremstyle{remark}
\newtheorem{remark}[theorem]{Remark}

\newcommand{\N}{\mathbb{N}}
\newcommand{\Z}{\mathbb{Z}}
\newcommand{\lean}[1]{\nolinkurl{#1}}
\newcommand{\Nzero}{4\,912\,740}

\raggedbottom

\begin{document}

\title[CST with few local maxima]{On Terras's coefficient stopping time conjecture for Collatz sequences with few local maxima}

\author[O. J. Said Duran]{Omar Javier Said Duran}
\address{Independent Researcher, Calama, Chile}
\email{omarsaid24@gmail.com}
\urladdr{https://orcid.org/0009-0009-1418-2558}

\date{12 October 2026}

\begin{abstract}
Let $T(n)=n/2$ for even $n$ and $T(n)=(3n+1)/2$ for odd $n$. The coefficient stopping time $\tau(n)$ is the least $k\ge1$ with $3^{q}<2^{k}$, where $q$ is the number of odd terms among $n,T(n),\dots,T^{k-1}(n)$. Terras conjectured in 1976 that $T^{\tau(n)}(n)<n$ for all $n\ge2$. Following a suggestion recorded by Rozier and Terracol, we study this coefficient stopping time (CST) conjecture according to the number of runs of odd terms, equivalently of local maxima, in the first $\tau(n)$ steps. We prove that $T^{\tau(n)}(n)<n$ whenever $\tau(n)$ is finite and there are at most $24$ runs. Adding the verification up to $2.8\cdot10^{19}$ by Rozier and Terracol, the same holds with at most $81$ runs; adding instead our own computation up to $3\cdot10^{11}$, with at most $42$ runs. For every fixed number of runs we give an explicit finite region containing all possible counterexamples. The method transfers that of Simons and de Weger for $m$-cycles to the first crossing: congruences along the runs, the bound of Laurent, Mignotte and Nesterenko for linear forms in two logarithms, one-sided approximations of $\log_2 3$, exact certificates organized as trees of residue classes, and a weighting of the runs adapted from Hercher. The proofs are formalized in Lean~4 with Mathlib, with the external bound and the finite verifications as explicit hypotheses and only the standard axioms. Nothing is proved about the finiteness of $\tau(n)$.
\end{abstract}

\subjclass[2020]{Primary 11B83; Secondary 11J86, 11Y55, 68V20}
\keywords{Collatz problem, $3x+1$ problem, stopping time, coefficient stopping time, linear forms in logarithms, formal verification, Lean}

\maketitle

\section{Introduction}\label{sec:intro}

Let $T$ be the Collatz map in its shortcut form, $T(n)=n/2$ for even $n$ and $T(n)=(3n+1)/2$ for odd $n$. For $n\ge1$ and $j\ge0$ let $S_j(n)$ be the number of odd terms among $n,T(n),\dots,T^{j-1}(n)$. Terras \cite{Ter76} introduced the \emph{stopping time} and the \emph{coefficient stopping time}
\[
\sigma(n)=\min\{k\ge1: T^{k}(n)<n\},\qquad \tau(n)=\min\{k\ge1: 3^{S_k(n)}<2^{k}\},
\]
with $\min\emptyset=\infty$. The name of $\tau$ comes from the affine formula $2^{j}T^{j}(n)=3^{S_j(n)}n+\beta_j(n)$ with $\beta_j(n)\ge0$ (Lemma~\ref{lem:affine}): the coefficient $3^{S_j(n)}/2^{j}$ predicts whether $T^{j}(n)$ lies below $n$, up to the additive term. Since $\beta_j(n)\ge0$, the orbit cannot descend below $n$ before the coefficient does, so $\tau(n)\le\sigma(n)$. Terras \cite{Ter76} and Everett \cite{Eve77} proved that $\sigma(n)<\infty$ for a set of $n$ of natural density one, and the Collatz conjecture is equivalent to $\sigma(n)<\infty$ for every $n\ge2$ \cite{Lag85,Lag10}.

Terras conjectured that $\sigma(n)=\tau(n)$ for all $n\ge2$ \cite{Ter76}. Lagarias \cite{Lag85} named this the \emph{coefficient stopping time} (CST) conjecture and noted that it excludes nontrivial cycles: the least element of a cycle other than $(1,2)$ has a finite coefficient stopping time but never descends (Corollary~\ref{cor:cycles} below). Rozier and Terracol \cite{RT26} call a finite sequence $n,T(n),\dots,T^{j}(n)$ \emph{paradoxical} if $3^{S_j(n)}<2^{j}$ and $T^{j}(n)\ge n$; a counterexample $n$ to the CST conjecture with $\tau(n)<\infty$ gives the paradoxical sequence of length $j=\tau(n)$. They verified the conjecture for $2\le n\le 2.8\cdot10^{19}$ \cite[Corollary~5.4]{RT26}, showed by an elementary argument that no acyclic paradoxical sequence has a single local maximum \cite[Appendix~A]{RT26}, and proposed, following a suggestion of Simons, to study the conjecture for sequences with a fixed number of local maxima \cite[Remark~5.5]{RT26}.

That approach has a long history for cycles. A cycle of $T$ is an \emph{$m$-cycle} if its parity vector has $m$ runs of odd terms, equivalently $m$ local maxima. Steiner \cite{Ste77} excluded nontrivial $1$-cycles using Baker's theory of linear forms in logarithms. Simons and de Weger \cite{SdW05} combined a lower bound for linear forms in two logarithms with congruences along the runs and a large computation, and excluded nontrivial $m$-cycles for $m\le68$ ($m\le75$ in a later version of their preprint); Hercher \cite{Her23} reached $m\le91$.

\subsection*{Results}
For $n\ge2$ with $\tau(n)=k<\infty$, let $r(n)$ be the number of runs of odd terms in the parity vector $(T^{i}(n)\bmod2)_{0\le i<k}$, that is, the number of maximal blocks of consecutive odd terms among $n,T(n),\dots,T^{k-1}(n)$. Since $T(x)>x$ for odd $x$ and $T(x)<x$ for even $x\ge2$, $r(n)$ is also the number of local maxima of the finite sequence $n,T(n),\dots,T^{k}(n)$ (Remark~\ref{rem:maxima}).

\begin{mainthm}\label{thm:A}
Let $n\ge2$ be an integer with $\tau(n)<\infty$ and $r(n)\le24$. Then $T^{\tau(n)}(n)<n$, that is, $\sigma(n)=\tau(n)$.
\end{mainthm}

\begin{mainthm}\label{thm:B}
Using in addition the verification of the CST conjecture for $2\le n\le 2.8\cdot10^{19}$ \cite[Corollary~5.4]{RT26}, the conclusion of Theorem~\ref{thm:A} holds for every $n\ge2$ with $\tau(n)<\infty$ and $r(n)\le81$.
\end{mainthm}

\begin{mainthm}\label{thm:D}
Suppose that $T^{\tau(n)}(n)<n$ for every $n$ with $2\le n\le3\cdot10^{11}$ and $\tau(n)<\infty$. Then the conclusion of Theorem~\ref{thm:A} holds for every $n\ge2$ with $\tau(n)<\infty$ and $r(n)\le42$.
\end{mainthm}

The hypothesis of Theorem~\ref{thm:D} is of course a consequence of \cite[Corollary~5.4]{RT26}; we verified it by a separate computation (Section~\ref{sec:census}). The point of Theorem~\ref{thm:D} is that it depends only on \cite{LMN95} and on a computation that can be repeated in a few hours on a personal computer.

Equivalently, a counterexample to the CST conjecture with finite coefficient stopping time has at least $25$ local maxima before its coefficient stopping time, at least $43$ given the computation of Section~\ref{sec:census}, and at least $82$ if \cite[Corollary~5.4]{RT26} is taken into account. All these theorems, and the next one, rest on the lower bound of Laurent, Mignotte and Nesterenko \cite{LMN95} (Theorem~\ref{thm:lmn}). For each fixed number of runs, that bound alone confines the possible counterexamples to an explicit finite region:

\begin{mainthm}\label{thm:C}
Let $n\ge2$ with $\tau(n)=k<\infty$ and $T^{k}(n)\ge n$, and let $r=r(n)$. Then
\[
k<2^{2r+32}\qquad\text{and}\qquad n<2^{19(2r+65)^2+r}.
\]
\end{mainthm}

Theorem~\ref{thm:C} is far from optimal, and Proposition~\ref{prop:tail} sharpens it; the point is that it holds for every $r$ with one proof. The cases $r\le24$, $r\le42$ and $r\le81$ are closed by exact computations inside the corresponding regions. As for cycles:

\begin{maincor}\label{cor:cycles}
Every cycle of $T$ other than $(1,2)$ has at least $25$ local maxima, and at least $82$ if \cite[Corollary~5.4]{RT26} is used.
\end{maincor}

Corollary~\ref{cor:cycles} is weaker than the results of \cite{SdW05,Her23} and is stated only as a consistency check: our method is theirs, applied to the first crossing instead of a full period.

\subsection*{What is not proved}
Theorems~\ref{thm:A}, \ref{thm:B} and~\ref{thm:D} assume that $\tau(n)$ is finite. A number with $\tau(n)=\infty$ would have an unbounded orbit (Remark~\ref{rem:infinite}); the finiteness of $\tau(n)$ for every $n$ is open, although it holds for a set of density one \cite{Ter76,Eve77}. The CST conjecture itself remains open for sequences with $25$ or more local maxima, and the theorems say nothing about whether orbits reach $1$.

\subsection*{Method}
The proof follows \cite{SdW05}. A failure of the CST conjecture at $k=\tau(n)$ forces $\Lambda n<r$, where $\Lambda=k\log2-q\log3>0$ and $q=S_k(n)$ (Corollary~\ref{cor:failure}). Each run of odd terms imposes a congruence on the value at its start, and these congruences bound $q$ in terms of $\log(n+1)$ and $r$ (Proposition~\ref{prop:sieve}). For large $k$ the lower bound of \cite{LMN95} makes $\Lambda$ too large for the two bounds to be compatible (Section~\ref{sec:region}). For moderate $q$ a one-sided best approximation lemma bounds $\Lambda$ from below on explicit intervals of $q$ (Section~\ref{sec:moderate}). What remains is a finite set of seeds, which is checked by certificates organized as trees of residue classes, with the congruence sieve restarted after each even term (Section~\ref{sec:certificates}). The tail and the table of intervals depend on $r$ only through explicit constants, while the size of the finite certificates grows quickly with $r$ (Remark~\ref{rem:cost}). For Theorems~\ref{thm:B} and~\ref{thm:D} the certificates are replaced by a weighting of the runs, adapted from an idea of \cite{Her23}: a run that starts far above $n$ contributes little to the failure inequality, so a failure needs many runs that start below a fixed power of two, and the sieve restarts at each of them (Section~\ref{sec:weights}).

\subsection*{Formal verification}
All the steps above, including the evaluation of the finite certificates, are formalized in Lean~4 with Mathlib \cite{Lean4,Mathlib}. The only external mathematical inputs are the specialization of the corollary of \cite{LMN95} to the logarithms of $2$ and $3$ and the finite verifications: \cite[Corollary~5.4]{RT26} for Theorem~\ref{thm:B}, and our computation for Theorem~\ref{thm:D}. All of them are explicit hypotheses of the formal statements. Section~\ref{sec:formal} gives the statements, the axioms used and the size of the computation.

\subsection*{Novelty}
Apart from the case of one local maximum in \cite[Appendix~A]{RT26}, we have not found a treatment of the CST conjecture by number of local maxima in the literature. Our search (arXiv, zbMATH Open and general search engines, October 2026) did not include MathSciNet. The weighting of Section~\ref{sec:weights} transfers one idea of \cite{Her23}; his merging lemma, which transmits convergence to $1$ but not the stopping time, does not seem to transfer, and the bounds of \cite{Her23} for cycles remain stronger than ours.

\section{The first crossing and the affine identity}\label{sec:affine}

Throughout, $\log$ is the natural logarithm and $\log_2$ the binary one. We write $\varepsilon_i(n)=T^{i}(n)\bmod 2$.

\begin{lemma}[Terras \cite{Ter76}]\label{lem:affine}
Let $n\ge1$ and $j\ge0$.
\begin{enumerate}
\item There is an integer $\beta_j(n)\ge0$ with $2^{j}T^{j}(n)=3^{S_j(n)}n+\beta_j(n)$.
\item If $0\le \rho<2^{j}$ and $u\ge0$, then $\varepsilon_i(\rho+2^{j}u)=\varepsilon_i(\rho)$ for $i<j$, $S_j(\rho+2^{j}u)=S_j(\rho)$ and $T^{j}(\rho+2^{j}u)=3^{S_j(\rho)}u+T^{j}(\rho)$.
\end{enumerate}
\end{lemma}

\begin{proof}
(1) Put $\beta_0=0$, and $\beta_{j+1}=3\beta_j+2^{j}$ if $T^{j}(n)$ is odd, $\beta_{j+1}=\beta_j$ otherwise; induction on $j$. (2) By induction on $i\le j$ one has $T^{i}(\rho+2^{j}u)=3^{S_i(\rho)}2^{j-i}u+T^{i}(\rho)$; for $i<j$ the first summand is even, so the parities agree, and both cases of $T$ preserve the form.
\end{proof}

\begin{lemma}\label{lem:before}
Let $n\ge2$ with $k=\tau(n)<\infty$.
\begin{enumerate}
\item $T^{j}(n)\ge n$ for $0\le j<k$.
\item If $n$ is even, then $k=1$ and $T(n)<n$.
\item If $n$ is odd, then $k\ge2$, $T^{k-1}(n)$ is even, and $q:=S_k(n)=S_{k-1}(n)$ satisfies $2^{k-1}\le3^{q}<2^{k}$. Hence $q\ge1$, $q<k$, and $\Lambda:=k\log2-q\log3$ satisfies $0<\Lambda\le\log2$.
\end{enumerate}
\end{lemma}

\begin{proof}
(1) For $1\le j<k$ we have $3^{S_j(n)}\ge2^{j}$, so $T^{j}(n)\ge 3^{S_j(n)}n/2^{j}\ge n$ by Lemma~\ref{lem:affine}. (2) $S_1(n)=0$ and $3^0<2$. (3) $S_1(n)=1$ and $3\ge2$, so $k\ge2$. If $T^{k-1}(n)$ were odd, then $3^{S_k(n)}=3\cdot3^{S_{k-1}(n)}\ge3\cdot2^{k-1}>2^{k}$. Thus $S_k(n)=S_{k-1}(n)$ and $2^{k-1}\le3^{q}<2^{k}$; the remaining claims follow.
\end{proof}

By (1), the CST conjecture says exactly that $T^{\tau(n)}(n)<n$ whenever $\tau(n)<\infty$. We call $n\ge2$ a \emph{failure} if $k=\tau(n)<\infty$ and $T^{k}(n)\ge n$. By (2), failures are odd.

\begin{definition}\label{def:runs}
Let $n\ge2$ be odd with $k=\tau(n)<\infty$. By Lemma~\ref{lem:before} the parity word $\varepsilon_0(n)\cdots\varepsilon_{k-1}(n)$ begins with $1$ and ends with $0$, so it can be written uniquely as
\[
1^{a_1}0^{b_1}1^{a_2}0^{b_2}\cdots1^{a_r}0^{b_r},\qquad a_i,b_i\ge1,
\]
and $r=r(n)\ge1$ is its number of \emph{runs}. For even $n$ we put $r(n)=0$. We write $A_0=0$, $A_i=a_1+\dots+a_i$ (so $A_r=q$), $\ell_1=0$, $\ell_{i+1}=\ell_i+a_i+b_i$ (so $\ell_{r+1}=k$), and $x_i=T^{\ell_i}(n)$ for $1\le i\le r+1$. In the formalization, $r(n)$ is computed directly from the orbit as the number of $i<k$ with $T^{i}(n)$ odd and either $i=0$ or $T^{i-1}(n)$ even.
\end{definition}

\begin{remark}\label{rem:maxima}
For $x\ge1$ odd, $T(x)=(3x+1)/2>x$, and for $x\ge2$ even, $T(x)<x$. Hence the sequence $n,T(n),\dots,T^{k}(n)$ increases exactly on the steps from odd terms. Its local maxima are the terms $T^{i}(n)$, $0<i<k$, with $T^{i-1}(n)$ odd and $T^{i}(n)$ even; the endpoints are not maxima because the word begins with $1$ and ends with $0$. So $r(n)$ is the number of local maxima.
\end{remark}

\begin{lemma}\label{lem:additive}
Let $n\ge2$ be odd with $k=\tau(n)<\infty$, $q=S_k(n)$ and $r=r(n)$. Then $2^{k}T^{k}(n)=3^{q}n+B$ with
\[
B=\sum_{i=1}^{r}2^{\ell_i}\,3^{q-A_i}\,\bigl(3^{a_i}-2^{a_i}\bigr),\qquad 0<B<r\,3^{q}.
\]
\end{lemma}

\begin{proof}
For odd $y$ we have $T(y)+1=\tfrac32(y+1)$, so a block of $a$ odd terms starting at $y$ gives $2^{a}(T^{a}(y)+1)=3^{a}(y+1)$, and a block of $b$ even terms divides by $2^{b}$. Hence
\[
2^{a_i+b_i}x_{i+1}=3^{a_i}x_i+3^{a_i}-2^{a_i}\qquad(1\le i\le r),
\]
and multiplying by $2^{\ell_i}$ and iterating gives the identity. Next,
\[
\frac{B}{3^{q}}=\sum_{i=1}^{r}\frac{2^{\ell_i}}{3^{A_{i-1}}}\Bigl(1-\bigl(\tfrac23\bigr)^{a_i}\Bigr).
\]
For $i=1$ the prefactor is $1$; for $i\ge2$ we have $1\le\ell_i<k$ and $S_{\ell_i}(n)=A_{i-1}$, so $2^{\ell_i}\le3^{A_{i-1}}$ by the definition of $\tau(n)$. Each summand lies in $(0,1)$.
\end{proof}

\begin{corollary}\label{cor:failure}
Let $n$ be a failure, $k=\tau(n)$, $q=S_k(n)$, $r=r(n)$, $\Delta=2^{k}-3^{q}$ and $\Lambda=k\log2-q\log3$. Then
\[
\Delta\, n\le B<r\,3^{q}\qquad\text{and}\qquad \Lambda\, n<r .
\]
\end{corollary}

\begin{proof}
$2^{k}n\le2^{k}T^{k}(n)=3^{q}n+B$ gives the first inequality, and $\Lambda\le e^{\Lambda}-1=\Delta/3^{q}$ gives the second.
\end{proof}

\begin{example}\label{ex:39}
The margin can be very small. For $n=39$ the orbit is $39,59,89,134,67,101,152,76,38$, with parity word $1110\,1100$; so $\tau(39)=8$, $q=5$, $3^{5}=243<256=2^{8}$, $r=2$, and $T^{8}(39)=38=n-1$.
\end{example}

\section{Congruences along the runs}\label{sec:sieve}

We keep the notation of Definition~\ref{def:runs}.

\begin{lemma}\label{lem:congruence}
Let $n\ge2$ be odd with $\tau(n)<\infty$. For $1\le i\le r$,
\[
2^{a_i}\mid x_i+1\qquad\text{and}\qquad 2^{a_i}(x_{i+1}+1)\le3^{a_i}(x_i+1).
\]
Consequently $2^{A_i}\le(n+1)\,3^{A_{i-1}}$ for $1\le i\le r$.
\end{lemma}

\begin{proof}
From $2^{a_i}(T^{a_i}(x_i)+1)=3^{a_i}(x_i+1)$ and $\gcd(2,3)=1$ we get $2^{a_i}\mid x_i+1$, and since $x_{i+1}=T^{a_i}(x_i)/2^{b_i}\le T^{a_i}(x_i)$, also the inequality. Chaining the inequalities gives $x_i+1\le(3/2)^{A_{i-1}}(n+1)$, and $x_i+1\ge2^{a_i}$ by divisibility; multiply by $2^{A_{i-1}}$.
\end{proof}

The congruences do not use the definition of $\tau(n)$; they hold for any word of the form $1^{a_1}0^{b_1}\cdots1^{a_r}0^{b_r}$ read along an orbit, which is used again in Section~\ref{sec:certificates}.

\begin{proposition}[the sieve]\label{prop:sieve}
Let $\alpha\ge1$ and $\beta\ge0$ be real numbers with $\log3\le\alpha\log2$, and define $C_0=0$ and $C_{j+1}=\lfloor\beta+\alpha C_j\rfloor$. If $n$ is odd with $\tau(n)<\infty$ and $\log(n+1)\le\beta\log2$, then $A_i\le C_i$ for $0\le i\le r$. The sequence $(C_j)$ is nondecreasing, so $q\le C_R$ whenever $r\le R$.
\end{proposition}

\begin{proof}
Taking $\log_2$ in Lemma~\ref{lem:congruence}, $A_{i}\le\log_2(n+1)+A_{i-1}\log_2 3\le\beta+\alpha C_{i-1}$, and $A_i$ is an integer. Monotonicity follows by induction from $C_1=\lfloor\beta\rfloor\ge0$.
\end{proof}

We use two instances. In the \emph{refined sieve}, $\alpha=317/200$ (note $\log_2 3<1.58497<1.585$) and $\beta=B/256$, where $B$ is an integer with $s^{256}\le2^{B}$ for some integer $s\ge n+1$; then
\begin{equation}\label{eq:refined}
C_{j+1}=\Bigl\lfloor\frac{25B+10144\,C_j}{6400}\Bigr\rfloor,
\end{equation}
and we write $C_R(B)$ for $C_R$. In the \emph{binary sieve} one uses $3^{5}<2^{8}$ and an integer $m$ with $n+1<2^{m}$: raising Lemma~\ref{lem:congruence} to the fifth power gives $5A_i<5m+8A_{i-1}$, hence $A_i\le c(i,m)$, where
\begin{equation}\label{eq:binary}
c(0,m)=0,\qquad c(j+1,m)=\Bigl\lfloor\frac{5m+8\,c(j,m)-1}{5}\Bigr\rfloor .
\end{equation}

\begin{lemma}\label{lem:geometric}
For all $j,m\ge0$, $3\cdot5^{j}c(j,m)+5^{j+1}m\le5\cdot8^{j}m$, and $c(j,m)+m\le2^{j}m$. In particular, if $5\cdot8^{R}\le(3A+5)5^{R}$, then $c(j,m)\le Am$ for all $j\le R$.
\end{lemma}

\begin{proof}
For the first inequality, the case $j=0$ is an equality. From $5c(j+1,m)\le5m+8c(j,m)$, multiplying by $3\cdot5^{j}$ and adding $5^{j+2}m$ gives eight times the left-hand side for $j$. The second follows from $c(j+1,m)+m\le2m+\tfrac85c(j,m)\le2(c(j,m)+m)$. The last claim uses monotonicity in $j$.
\end{proof}

\begin{lemma}\label{lem:time}
If $q\ge1$ and $2^{k-1}\le3^{q}$, then $5k\le8q+5$; in particular $k\le2q+1$.
\end{lemma}

\begin{proof}
$2^{5(k-1)}\le3^{5q}=243^{q}<256^{q}=2^{8q}$.
\end{proof}

\section{Linear forms in two logarithms}\label{sec:lmn}

The following is Corollary~2 of \cite{LMN95}, which we checked against the original (p.~288).

\begin{theorem}[Laurent, Mignotte, Nesterenko]\label{thm:lmn}
Let $\alpha_1,\alpha_2$ be multiplicatively independent positive real algebraic numbers, $b_1,b_2$ positive integers and $\Lambda=b_2\log\alpha_2-b_1\log\alpha_1$. Put $D=[\mathbb{Q}(\alpha_1,\alpha_2):\mathbb{Q}]/[\mathbb{R}(\alpha_1,\alpha_2):\mathbb{R}]$ and let $A_1,A_2>1$ satisfy $\log A_i\ge\max\{h(\alpha_i),|\log\alpha_i|/D,1/D\}$, where $h$ is the absolute logarithmic height. With $b'=b_1/(D\log A_2)+b_2/(D\log A_1)$,
\[
\log|\Lambda|\ge-24.34\,D^{4}\bigl(\max\{\log b'+0.14,\,21/D,\,1/2\}\bigr)^{2}\log A_1\log A_2 .
\]
\end{theorem}

\begin{lemma}\label{lem:lmn23}
For all integers $k,q\ge1$,
\[
\log|k\log2-q\log3|\ge-24.34\,\log3\,\bigl(\max\{\log(q+k/\log3)+0.14,\,21\}\bigr)^{2}.
\]
\end{lemma}

\begin{proof}
Take $\alpha_1=3$, $\alpha_2=2$, $b_1=q$, $b_2=k$ in Theorem~\ref{thm:lmn}. Then $D=1$, $h(3)=\log3>1$ and $h(2)=\log2<1$, so we may take $A_1=3$ and $A_2=e$. The numbers $2$ and $3$ are multiplicatively independent, $b'=q+k/\log3$, and $\max\{\cdot,21,1/2\}=\max\{\cdot,21\}$.
\end{proof}

\begin{lemma}\label{lem:weak}
Let $k\ge2$ and $q$ be integers with $3^{q}<2^{k}\le2\cdot3^{q}$, put $\Lambda=k\log2-q\log3$ and $t=\lfloor\log_2k\rfloor$. Then
\[
\Lambda\ge\exp\bigl(-27(\log(2k)+21)^{2}\bigr)\ge2^{-19(t+33)^{2}}.
\]
\end{lemma}

\begin{proof}
The hypotheses give $1\le q<k$ and $\Lambda>0$. Since $\log3>1$, $b'=q+k/\log3<2k$, so the maximum in Lemma~\ref{lem:lmn23} is at most $\log(2k)+21$; and $24.34\log3<26.75\le27$. This gives the first inequality. For the second, $\log(2k)\le(t+2)\log2$ and $21\le31\log2$, so $\log(2k)+21\le(t+33)\log2$, and $27\log2<18.72\le19$.
\end{proof}

\begin{corollary}\label{cor:lmnfail}
A failure $n$ with $k=\tau(n)$, $r=r(n)$ and $t=\lfloor\log_2k\rfloor$ satisfies $n<r\cdot2^{19(t+33)^{2}}$.
\end{corollary}

\begin{proof}
By Lemma~\ref{lem:before}(3), $k\ge2$ and $3^{q}<2^{k}\le2\cdot3^{q}$; combine Corollary~\ref{cor:failure} and Lemma~\ref{lem:weak}.
\end{proof}

\section{A finite region for each number of runs}\label{sec:region}

\begin{proof}[Proof of Theorem~\ref{thm:C}]
Let $n$ be a failure with $k$, $q$, $r$, $t$ as above, and put $E=19(t+33)^{2}$. By Corollary~\ref{cor:lmnfail}, $n<r2^{E}\le2^{E+r}$, so $n+1<2^{E+r+1}$ and, by the binary sieve and Lemma~\ref{lem:geometric}, $q\le c(r,E+r+1)<2^{r}(E+r+1)$. By Lemma~\ref{lem:time}, $2^{t}\le k\le2q+1$. Suppose $t\ge2r+32$. Then $r\le t$ and $E+r+1\le20(t+33)^{2}$, so
\[
2^{t}\le40\cdot2^{r}(t+33)^{2}+1 .
\]
But $40\cdot2^{r}(t+33)^{2}+1<2^{t}$ for all $t\ge2r+32$: at $r=0$, $t=32$ the margin is $2^{32}-40\cdot65^{2}-1=4\,294\,798\,295$; raising $r$ by one multiplies the right-hand side at $t=2r+32$ by $4$ and the left by at most $4$, because $(2r+67)^{2}\le2(2r+65)^{2}$; and raising $t$ by one doubles the right-hand side, while $(t+34)^{2}\le2(t+33)^{2}$. Hence $t\le2r+31$, that is, $k<2^{2r+32}$, and then $n<2^{E+r}\le2^{19(2r+65)^{2}+r}$.
\end{proof}

The following criterion replaces the crude constants by a certificate of three integer inequalities, and is what we use for $R=24$, $R=42$ and $R=81$.

\begin{proposition}\label{prop:tail}
Let $R,A,c,H$ be natural numbers with
\[
R\le2^{c},\qquad 5\cdot8^{R}\le(3A+5)\,5^{R},\qquad 8A\bigl(19(H+33)^{2}+c+1\bigr)+5<5\cdot2^{H}.
\]
If $n$ is a failure with $r(n)\le R$, then $\tau(n)<2^{H}$ and $n<2^{19(H+32)^{2}+c}$.
\end{proposition}

\begin{proof}
With $E=19(t+33)^{2}$, Corollary~\ref{cor:lmnfail} gives $n<R\,2^{E}\le2^{E+c}$, so $n+1<2^{E+c+1}$. By the binary sieve, monotonicity in the number of runs and Lemma~\ref{lem:geometric}, $q\le c(R,E+c+1)\le A(E+c+1)$, and by Lemma~\ref{lem:time},
\[
5\cdot2^{t}\le5k\le8A\bigl(19(t+33)^{2}+c+1\bigr)+5 .
\]
The third hypothesis contradicts this for $t=H$. The contradiction persists for all $t\ge H$, because doubling the left-hand side dominates the change on the right:
\[
2\bigl(19(t+33)^{2}+c+1\bigr)-\bigl(19(t+34)^{2}+c+1\bigr)=19t^{2}+1216t+19419+c\ge0 .
\]
Hence $t<H$, so $k<2^{H}$ and $t+33\le H+32$.
\end{proof}

For $R=24$ the constants $(A,c,H)=(132\,046,\,5,\,35)$ satisfy the hypotheses, for $R=42$ the constants $(623\,574\,031,\,7,\,51)$, and for $R=81$ the constants $(56\,959\,654\,291\,224\,268,\,7,\,75)$. In general one may take $A=\lceil5(8^{R}-5^{R})/(3\cdot5^{R})\rceil$ and the least $c$ with $R\le2^{c}$; the least admissible $H$ then grows roughly linearly in $R$ (for instance $H=88$ for $R=100$ and $H=703$ for $R=1000$), much more slowly than the exponent $2R+32$ of Theorem~\ref{thm:C}.

\section{Moderate values of \texorpdfstring{$q$}{q}: one-sided approximations}\label{sec:moderate}

By Proposition~\ref{prop:tail} a failure with $r(n)\le R$ has $q<k<2^{H}$. To bound $n$ for $q$ below this limit we use that $\Lambda$ is \emph{positive} at the first crossing.

\begin{lemma}\label{lem:onesided}
Let $x,y$ be real, $L(a,b)=ax-by$, and let $(P,Q)$ and $(p,q')$ be integer pairs with $Q,q'>0$, $Pq'-pQ=1$ and $L(P,Q)>0>L(p,q')$. If $(k,b)$ are integers with $0<b<Q+q'$ and $L(k,b)>0$, then $L(k,b)\ge L(P,Q)$.
\end{lemma}

\begin{proof}
Write $(k,b)=u(P,Q)+v(p,q')$ with $u=kq'-pb$ and $v=Pb-kQ$, which is possible because the determinant is $1$. If $u\le0$, then $b>0$ forces $v\ge1$, and $L(k,b)=uL(P,Q)+vL(p,q')<0$, a contradiction. So $u\ge1$. If $v\ge1$, then $b\ge Q+q'$, again a contradiction. So $v\le0$ and $L(k,b)\ge uL(P,Q)\ge L(P,Q)$.
\end{proof}

No coprimality of $k$ and $b$ is needed, and the pairs need not be convergents. With $x=\log2$, $y=\log3$ and $(k,b)=(\tau(n),q)$, Lemma~\ref{lem:onesided} bounds $\Lambda$ from below on the interval $Q\le q<Q+q'$.

\begin{proposition}\label{prop:row}
Let $(P,Q)$, $(p,q')$ be as in Lemma~\ref{lem:onesided} for $x=\log2$, $y=\log3$, let $c_0>0$ be a rational number with $c_0\le P\log2-Q\log3$, let $R\ge1$, $s=\lceil R/c_0\rceil$, and let $B$ be an integer with $s^{256}\le2^{B}$. If $n$ is a failure with $r(n)\le R$ and $Q\le q<Q+q'$, then $n<s$ and $q\le C_R(B)$. In particular, if $C_R(B)<Q$ there is no such failure.
\end{proposition}

\begin{proof}
By Lemma~\ref{lem:onesided}, $\Lambda\ge c_0$, so $n<r/\Lambda\le R/c_0\le s$ by Corollary~\ref{cor:failure}. Then $n+1\le s$ and the refined sieve applies with $\beta=B/256$.
\end{proof}

The signs of $L(P,Q)$ and $L(p,q')$ and the rational numbers $c_0$ are certified with enclosures of $\log2$ and $\log3$ to $55$ decimal places, which are proved in Lean from the power series of $\log(1-x)$ with Mathlib's explicit remainder bound. The pairs are produced by a mediant (Stern--Brocot) process for $\log_23$, and consecutive rows tile an interval of values of $q$.

For $R=24$ the table has $43$ rows covering $1\le q<39\,123\,338\,641$; it is reproduced in Appendix~\ref{app:table}. In $35$ rows the bound is $s\le N_0+1$, where $N_0=\Nzero$; in $6$ rows $C_{24}(B)<Q$; and two rows leave larger seeds: $79\,335\le q<190\,537$ with $n\le6\,548\,930$, and $190\,537\le q<10\,781\,274$ with $n\le372\,049\,745$. For $R=81$ the table has $101$ rows and covers $1\le q<Q^{*}$, where $Q^{*}=38\,887\,293\,218\,019\,448\,937\,881$. In $91$ rows the bound is $n\le2.8\cdot10^{19}$; the other $10$ rows, which start at $q=130\,441\,933\,147\,714\,940$, are closed in Section~\ref{sec:weights}. For $R=42$ the table has $64$ rows and covers $1\le q<2\,365\,627\,332\,311\,699$; in $41$ rows the bound is $n\le3\cdot10^{11}$, and the other $23$ rows, from $q=397\,573\,379$, are also closed in Section~\ref{sec:weights}. Since $39\,123\,338\,641>2^{35}$, $Q^{*}>2^{75}$ and $2\,365\,627\,332\,311\,699>2^{51}$, each table reaches the range of Proposition~\ref{prop:tail}.

\section{Finite certificates}\label{sec:certificates}

\subsection{Residue classes}
By Lemma~\ref{lem:affine}(2), the integers $n\ge\rho$ with $n\equiv\rho\pmod{2^{i}}$ share their first $i$ parities, and $T^{i}(n)=au+c$ for $n=\rho+mu$, $u\ge0$, where $m=2^{i}$, $a=3^{S_i(\rho)}$ and $c=T^{i}(\rho)$. A \emph{state} $(\rho,m,a,c)$ represents this class. Writing $u=2v+b$ with $b\in\{0,1\}$ splits it into the two states with residues $\rho$ and $\rho+m$, modulus $2m$, and, according to the parity of $ab+c$, multiplier $a$ or $3a$ and constant $T(ab+c)$. Along a branch on which $a\ge m$ at every earlier node, the first node with $a<m$ is the first crossing for every member of the class, and if moreover $c<\rho$, then $au+c<\rho+mu$ for all $u\ge0$: the whole class satisfies $T^{\tau(n)}(n)<n$.

\subsection{Small seeds}

\begin{proposition}\label{prop:small}
For $2\le n\le N_0=\Nzero$, $\tau(n)\le224$ and $T^{\tau(n)}(n)<n$.
\end{proposition}

This is a finite statement. It is proved in Lean by a verified tree of states starting from $(2,1,1,2)$: a node is accepted if its class contains no integer up to $N_0$, or if the coefficient has crossed and $c<\rho$; when only one integer of the class remains below $N_0$, it is checked individually by iteration; a node whose coefficient has crossed with $c\ge\rho$, or whose fuel is exhausted, is rejected. A theorem proved in Lean shows that acceptance of the root implies the proposition. The tree has $243\,225$ nodes, with $15\,466$ shared leaves and $106\,147$ individual leaves. No bound for linear forms in logarithms is used here, and the number of runs plays no role. An independent program reconstructs the $121\,613$ classes, checks that they are disjoint and cover exactly $4\,912\,739$ integers, and iterates two representatives of each.

\subsection{Restarting the sieve after an even term}
Let $n$ be odd with $k=\tau(n)<\infty$ and let $0\le i<k$ be a \emph{cut}: $i=0$, or $T^{i-1}(n)$ is even. Then no run crosses the cut, and
\[
r(n)=s_i+r_i'(n),\qquad s_i=\#\{\text{runs that start before } i\},
\]
where $r_i'(n)$ is the number of runs of the word $\varepsilon_i(n)\cdots\varepsilon_{k-1}(n)$. That word ends with $0$, so after removing its initial zeros it has the shape of Definition~\ref{def:runs}, and Lemma~\ref{lem:congruence} and Proposition~\ref{prop:sieve} apply to it with $T^{i}(n)$ in place of $n$; no first-crossing property of the shifted word is needed.

\begin{lemma}\label{lem:restart}
Let $i$ be a cut, $r(n)\le R$, and let $s,B$ be integers with $T^{i}(n)+1\le s$ and $s^{256}\le2^{B}$. Then $S_k(n)\le S_i(n)+C_{R-s_i}(B)$.
\end{lemma}

\subsection{The window verifier}
A \emph{window} is given by integers $N_1\le N_2$, a number of runs $R$ and a threshold $q_{\min}$; it asserts that there is no failure $n$ with $N_1\le n\le N_2$, $r(n)\le R$ and $S_{\tau(n)}(n)\ge q_{\min}$. The verifier explores states enriched with $S_i$, $s_i$ and the parity of the previous term, starting from the class of all $n\ge N_1$. For a state of index $i$, the largest value of $T^{i}(n)$ for members $n\le N_2$ is $X=a\lfloor(N_2-\rho)/m\rfloor+c$, and $B$ is the least integer with $(X+1)^{256}\le2^{B}$, computed by bisection. A node is
\begin{enumerate}
\item accepted if $\rho>N_2$ (no member in the window);
\item if $a<m$ (the first crossing), accepted when $c<\rho$ and rejected otherwise;
\item accepted if $s_i>R$;
\item accepted if $i$ is a cut and $S_i+C_{R-s_i}(B)<q_{\min}$;
\end{enumerate}
and otherwise split into its two children. Exhausting the fuel rejects.

\begin{proposition}\label{prop:window}
If the verifier accepts the root of a window, the window assertion holds.
\end{proposition}

\begin{proof}[Sketch]
By induction on the fuel: an accepted node cannot contain a failure with the stated properties, by the remark at the end of the first subsection for rule (2), by $r(n)\ge s_i$ for rule (3), and by Lemma~\ref{lem:restart} for rule (4), since $S_{\tau(n)}(n)\ge q_{\min}$; a split node is accepted only if both children are. The full proof is formalized (Section~\ref{sec:formal}).
\end{proof}

Since $B$ enters the recurrence~\eqref{eq:refined} monotonically, a certificate may also supply at each leaf any $B$ with $(X+1)^{256}\le2^{B}$, which the kernel checks directly. The windows needed for Theorem~\ref{thm:A} are listed in Table~\ref{tab:windows}. The windows $W_1$ and $W_2$ overlap; only the threshold $q_{\min}$ distinguishes them.

\begin{table}[ht]
\caption{Windows certified by the verifier. Seeds are counted inclusively; ``leaves'' are the accepted leaves that contain integers of the window.}\label{tab:windows}
\small
\begin{tabular}{@{}lrrrrrrr@{}}
\toprule
& $N_1$ & $N_2$ & $q_{\min}$ & $R$ & nodes & leaves & depth \\
\midrule
$W_1$ & $4\,912\,741$ & $6\,548\,930$ & $79\,335$ & $24$ & $756\,925$ & $47\,048$ & $62$ \\
$W_2$ & $4\,912\,741$ & $372\,049\,745$ & $190\,537$ & $24$ & $29\,403\,783$ & $3\,939\,127$ & $65$ \\
\bottomrule
\end{tabular}
\par\smallskip
\raggedright\footnotesize The fuel was $128$.
\end{table}

\subsection{Evaluation in the kernel}
The trees are cut into subtrees of at most $2\,000$ nodes. The acceptance of each subtree is a closed Boolean computation, proved with \texttt{decide +kernel}, so it is evaluated by the Lean kernel without compiled code; the acceptance of the root then follows from a proved lemma that combines the conclusions of two children. For $W_1$ and $W_2$ this gives $664$ and $25\,275$ subtrees, joined by $663$ and $25\,274$ unions, in $39$ and $596$ modules.

\section{Weights of runs}\label{sec:weights}

In this section $n$ is odd with $k=\tau(n)<\infty$, and we use the notation of Definition~\ref{def:runs}. Let $M$ and $R\ge r$ be positive integers. The $i$-th run is \emph{heavy} if $x_i+1\le M$ and \emph{light} otherwise; let $h$ be the number of heavy runs, and put $w_i=M$ if the $i$-th run is heavy and $w_i=n+R$ if it is light.

\begin{lemma}\label{lem:weights}
If $n$ is a failure with $r(n)\le R$, then
\[
\Delta\,n\,M<3^{q}\bigl(hM+(r-h)(n+R)\bigr).
\]
Consequently, if $c\le\Lambda$, then $c\,n\,M<hM+(r-h)(n+R)$.
\end{lemma}

\begin{proof}
Put $W_1=0$ and $W_{i+1}=W_i+w_i$. We show by induction that $2^{\ell_i}x_iM\le3^{A_{i-1}}(nM+W_i)$ for $1\le i\le r+1$, with strict inequality for $i\ge2$. The case $i=1$ is an equality. By the block identity in the proof of Lemma~\ref{lem:additive},
\[
2^{\ell_{i+1}}x_{i+1}M=3^{a_i}\,2^{\ell_i}x_iM+\bigl(3^{a_i}-2^{a_i}\bigr)2^{\ell_i}M<3^{a_i}\bigl(2^{\ell_i}x_iM+2^{\ell_i}M\bigr),
\]
so it suffices to show $2^{\ell_i}M\le3^{A_{i-1}}w_i$. For $i\ge2$ we have $2^{\ell_i}\le3^{A_{i-1}}$ because $\ell_i<k$, and for $i=1$ both sides of this inequality equal $1$. This settles heavy runs. If the run is light, then $M<x_i+1$; the partial sums of the expression for $B$ in Lemma~\ref{lem:additive} give $2^{\ell_i}x_i<3^{A_{i-1}}(n+i-1)$ for $i\ge2$, so $2^{\ell_i}(x_i+1)\le3^{A_{i-1}}(n+i)\le3^{A_{i-1}}(n+R)$, and this also holds for $i=1$. Hence $2^{\ell_i}M<3^{A_{i-1}}(n+R)$. At $i=r+1$ we get $2^{k}T^{k}(n)M<3^{q}(nM+W_{r+1})$ with $W_{r+1}=hM+(r-h)(n+R)$, and $T^{k}(n)\ge n$ gives the first inequality. The second follows from $\Lambda\le e^{\Lambda}-1=\Delta/3^{q}$.
\end{proof}

Each light run contributes less than $(n+R)/M$ to the right-hand side, so when $M$ is much larger than $n$, a failure with large $\Lambda n$ needs about $\Lambda n$ heavy runs. The next lemma restarts the sieve at each of them. For integers $m\ge1$ and $j\ge0$ let
\[
g(m,j)=j+C_j\bigl(256(m-1)\bigr),
\]
with $C_j(B)$ the refined sieve~\eqref{eq:refined}.

\begin{lemma}[discounted sieve]\label{lem:discounted}
Let $x\ge1$ be odd and suppose that the parity word of $x$ over some number of steps has the form $1^{a_1}0^{b_1}\cdots1^{a_j}0^{b_j}$ with $a_i,b_i\ge1$. If $x+1\le2^{m}$ with $m\ge1$, then the number of odd terms $a_1+\dots+a_j$ is at most $g(m,j)$.
\end{lemma}

\begin{proof}
With $x_1=x$ and $x_{i+1}$ the value after the $i$-th block, $x_{i+1}\le T^{a_i}(x_i)/2$ because $b_i\ge1$, and $2^{a_i}\mid x_i+1$ as in Lemma~\ref{lem:congruence}. From $2^{a}(T^{a}(y)+1)=3^{a}(y+1)$ one gets $2^{a_i}(x_{i+1}+1)\le\tfrac12\,3^{a_i}(x_i+1)+2^{a_i-1}\le2\cdot3^{a_i-1}(x_i+1)$, where the last step uses $2^{a_i}\le x_i+1$. So $2^{a_i-1}(x_{i+1}+1)\le3^{a_i-1}(x_i+1)$. Writing $D_i=\sum_{l\le i}(a_l-1)$ and chaining as in Lemma~\ref{lem:congruence}, $2^{D_i+1}\le(x+1)\,3^{D_{i-1}}$, so $D_i\le(m-1)+D_{i-1}\log_23$. Proposition~\ref{prop:sieve} with $\beta=m-1$ and $\alpha=317/200$ gives $D_j\le C_j(256(m-1))$, and the number of odd terms is $D_j+j$.
\end{proof}

The increments of $g(m,\cdot)$ are nondecreasing, and $g(m,1)=m$. Since a run starts after an even term, the word from the start of a run to the end has the form of Lemma~\ref{lem:discounted}; cutting the word of $n$ at the starts of the heavy runs and using convexity gives:

\begin{lemma}\label{lem:segments}
If $n+1\le M=2^{m}$ and at least $h_1\ge1$ runs are heavy, then $q\le g(m,r-h_1+1)+(h_1-1)m$.
\end{lemma}

\begin{proposition}[pieces]\label{prop:piece}
Let $(P,Q)$, $(p,q')$ and $c_0$ be as in Proposition~\ref{prop:row}, and let $R\ge1$, $1\le N_1\le N_2$, $m\ge1$ and $0\le h_0\le R$ be integers with $N_2+R\le2^{m}$,
\[
h_0\,2^{m}+(R-h_0)(N_2+R)\le c_0\,N_1\,2^{m},
\]
and, if $h_0<R$, also $g(m,R-h_0)+h_0m<Q$.
Then there is no failure $n$ with $N_1\le n\le N_2$, $r(n)\le R$ and $Q\le q<Q+q'$.
\end{proposition}

\begin{proof}
For such a failure $\Lambda\ge c_0$ by Lemma~\ref{lem:onesided}. Apply Lemma~\ref{lem:weights} with $M=2^{m}$. If $h\le h_0$, then, since $n+R\le M$ and $r\le R$, the right-hand side is at most $h_0M+(R-h_0)(N_2+R)\le c_0N_1M\le c_0nM$, a contradiction. So $h\ge h_0+1$, which forces $h_0<R$, and the run starting at $n$ is heavy. By Lemma~\ref{lem:segments} with $h_1=h_0+1$ and monotonicity of $g$, $q\le g(m,R-h_0)+h_0m<Q$, a contradiction.
\end{proof}

The certificates for Theorems~\ref{thm:B} and~\ref{thm:D} cut the seeds of each remaining row, from the external bound up to $\lceil R/c_0\rceil$, into consecutive pieces $[N_1,N_2]$ with their own $m$ and $h_0$. For $R=81$ the $10$ rows use $119$ pieces, and for $R=42$ the $23$ rows use $270$ pieces; in most pieces $h_0=0$, that is, a single segment from $n$ already suffices. The certificates are generated with exact rational arithmetic, and an independent program recomputes $\log2$ and $\log3$ by $\operatorname{artanh}$ series with bounded remainders, the determinants and signs, the contiguity and coverage of the pieces, the sieve and the constants of the tail.

\begin{remark}[The limit of the weights]\label{rem:limit}
For $R=82$ the method fails in a single row, $397\,560\,349\,370\,386\,783\le q<4\,640\,282\,259\,296\,926\,456$, on the seeds $2.8\cdot10^{19}<n\le32\,962\,697\,863\,578\,796\,386$. There $c_0n<5$, so five heavy runs suffice, and five heavy runs still allow $q\le4.43\cdot10^{17}$. A heavy run whose coefficient is close to $1$ weighs almost $1$, and in that row such coefficients occur, so refining the weights does not help; closing $82$ needs another tool in that window or a larger external range.
\end{remark}

\section{A computation up to \texorpdfstring{$3\cdot10^{11}$}{3e11}}\label{sec:census}

The hypothesis of Theorem~\ref{thm:D} was checked by a program written for this purpose, without using the data of \cite{RT26}. By Lemma~\ref{lem:affine}(2), all $n\equiv\rho\pmod{2^{24}}$ share their first $24$ parities. For the $16\,490\,635$ residues whose coefficient crosses at some step $j\le24$, every member $n=\rho+2^{24}u\ge2$ has $\tau(n)=j$ and $2^{j}T^{j}(n)=3^{S}n+\beta$ with $S$ and $\beta$ fixed, so it descends exactly when $n(2^{j}-3^{S})>\beta$; this is checked exactly for the least member, and it then holds for all larger ones. The members of the remaining $286\,581$ residue classes are iterated one by one until their first crossing, $5\,124\,467\,571$ seeds in total; the few iterations that would overflow $64$-bit integers ($1\,506$ in total) are redone with exact integers. No failure was found. The range was covered in thirteen pieces: $[2,2^{36}]$ in $42$ minutes, and $(2^{36},2\cdot10^{11}]$ and $(2\cdot10^{11},3\cdot10^{11}]$ in six parallel processes each, of about $31$ and $27$ minutes. A separate script checks that the pieces are contiguous and that each piece iterated exactly the members of the residue classes without crossing in its range.

To test that the program detects failures when they exist, it can be run with $T^{k}(n)<n$ replaced by $fT^{k}(n)<gn$; for $f/g=4/3$, $6/5$ and $10/9$ it found exactly the same numbers of failures as brute force. An independent program, which does not import the first one, iterates every $n\le2^{27}$ without any sieve and three million random seeds spread over the three pieces; it found no failure. This computation is not new: much larger ranges were verified in \cite{RT26}. It is used only to make Theorem~\ref{thm:D} independent of \cite{RT26}.

\section{Proofs of Theorems~\ref{thm:A}, \ref{thm:B} and~\ref{thm:D} and Corollary~\ref{cor:cycles}}\label{sec:proofs}

\begin{proof}[Proof of Theorem~\ref{thm:A}]
Suppose that $n$ is a failure with $r(n)\le24$, and let $k=\tau(n)$ and $q=S_k(n)$. Then $n$ is odd, and Proposition~\ref{prop:tail} with $(R,A,c,H)=(24,132\,046,5,35)$ gives $k<2^{35}$; since $q<k$, $q<39\,123\,338\,641$, so $q$ lies in one of the $43$ rows of Appendix~\ref{app:table}. Proposition~\ref{prop:row} excludes the six rows marked ``sieve'' and gives $n\le N_0$ in $35$ rows, contradicting Proposition~\ref{prop:small}. In the remaining two rows, either $n\le N_0$, again impossible, or $n$ lies in $W_1$ (if $q\ge79\,335$ and $n\le6\,548\,930$) or in $W_2$ (if $q\ge190\,537$ and $n\le372\,049\,745$), contradicting Proposition~\ref{prop:window}.
\end{proof}

\begin{proof}[Proof of Theorem~\ref{thm:B}]
Suppose that $n$ is a failure with $r(n)\le81$. Proposition~\ref{prop:tail} with $R=81$, $c=7$, $H=75$ and $A=56\,959\,654\,291\,224\,268$ gives $k<2^{75}$, so $q<Q^{*}$ and $q$ lies in one of the $101$ rows of Section~\ref{sec:moderate}. In $91$ rows Proposition~\ref{prop:row} gives $n\le2.8\cdot10^{19}$, which contradicts \cite[Corollary~5.4]{RT26}. In the other $10$ rows either $n\le2.8\cdot10^{19}$, with the same contradiction, or $n$ lies in one of the $119$ pieces, contradicting Proposition~\ref{prop:piece}.
\end{proof}

\begin{proof}[Proof of Theorem~\ref{thm:D}]
The same argument with $R=42$, $c=7$, $H=51$, $A=623\,574\,031$, the $64$ rows of Section~\ref{sec:moderate}, the hypothesis $n\le3\cdot10^{11}$ in place of \cite[Corollary~5.4]{RT26}, and $270$ pieces.
\end{proof}

\begin{proof}[Proof of Corollary~\ref{cor:cycles}]
Let $\mu$ be the least element of a cycle of length $L$ other than $(1,2)$, with $K$ odd terms and $M$ local maxima. Then $\mu$ is odd, and the term before $\mu$ in the cycle is even, since an odd predecessor $x$ would satisfy $x=(2\mu-1)/3<\mu$. So the parity word of the cycle read from $\mu$ begins with $1$, ends with $0$ and has $M$ runs. Lemma~\ref{lem:affine} over one period gives $2^{L}\mu=3^{K}\mu+\beta_L(\mu)$ with $\beta_L(\mu)>0$, so $3^{K}<2^{L}$ and $\tau(\mu)\le L$; the parity word up to $\tau(\mu)$ is a prefix of the cycle word, with at most $M$ runs. Since the orbit never goes below $\mu$, $\mu$ is a failure, and Theorems~\ref{thm:A} and~\ref{thm:B} give $M\ge25$, respectively $M\ge82$.
\end{proof}

\section{Formal verification}\label{sec:formal}

The proofs are formalized in Lean~4.32.2 (commit \lean{f3b06c705e6c}) with Mathlib at commit \lean{905b95818eb3}. The identifiers are in Spanish. In the statements below, \texttt{PrimerCruce n k} says that $\tau(n)=k$, and \texttt{rachas n k} is the number of runs of $n$ up to step $k$.

\begin{Verbatim}[fontsize=\small]
def T (n : ℕ) : ℕ := if n % 2 = 0 then n / 2 else (3 * n + 1) / 2
def S : ℕ → ℕ → ℕ
  | 0, _ => 0
  | k + 1, n => n % 2 + S k (T n)
def PrimerCruce (n k : ℕ) : Prop :=
  (∀ j, j < k → 1 ≤ j → 2 ^ j ≤ 3 ^ S j n) ∧ 3 ^ S k n < 2 ^ k
def rachas (n : ℕ) : ℕ → ℕ
  | 0 => 0
  | k + 1 => rachas n k +
      if T^[k] n % 2 = 1 ∧ (k = 0 ∨ T^[k - 1] n % 2 = 0) then 1 else 0
def LMN23 : Prop :=
  ∀ k q : ℕ, 0 < k → 0 < q →
    -(24.34 * (max (log ((q : ℝ) + k / log 3) + 0.14)
      (max 21 (1 / 2))) ^ 2 * log 3) ≤
        log |(k : ℝ) * log 2 - q * log 3|
def RozierTerracol54 : Prop :=
  ∀ {n k : ℕ}, 2 ≤ n → n ≤ 28000000000000000000 →
    PrimerCruce n k → T^[k] n < n
theorem cst_veinticuatro_corolario (hL : LMN23) {n k : ℕ}
    (hn : 2 ≤ n) (hk : PrimerCruce n k) (hr : rachas n k ≤ 24) :
    (∀ j, j < k → n ≤ T^[j] n) ∧ T^[k] n < n
theorem cst_condicional (hL : LMN23) (hRT : RozierTerracol54)
    {n k : ℕ} (hn : 2 ≤ n) (hk : PrimerCruce n k)
    (hr : rachas n k ≤ 81) :
    (∀ j, j < k → n ≤ T^[j] n) ∧ T^[k] n < n
def CensoPropio3e11 : Prop :=
  ∀ {n k : ℕ}, 2 ≤ n → n ≤ 300000000000 →
    PrimerCruce n k → T^[k] n < n
theorem cst_condicional (hL : LMN23) (hC : CensoPropio3e11)
    {n k : ℕ} (hn : 2 ≤ n) (hk : PrimerCruce n k)
    (hr : rachas n k ≤ 42) :
    (∀ j, j < k → n ≤ T^[j] n) ∧ T^[k] n < n
\end{Verbatim}

\noindent The first theorem is in the module \lean{OcelotasVeinticuatro}, together with \lean{fallo_exige_veinticinco}, which states that a failure has at least $25$ runs; the second is in \lean{OcelotasOchentaYUno}, together with \lean{fallo_exige_ochenta_y_dos}, and the third in \lean{OcelotasCuarentaYDos}, together with \lean{fallo_exige_cuarenta_y_tres}. Lemmas~\ref{lem:weights}--\ref{lem:segments} and Proposition~\ref{prop:piece} are in \lean{OcelotasPesos}. Line breaks are adjusted for display. \texttt{LMN23} is Lemma~\ref{lem:lmn23}, and \lean{debil_de_corolario} derives from it the first inequality of Lemma~\ref{lem:weak}; the passage from Theorem~\ref{thm:lmn} to Lemma~\ref{lem:lmn23} (Section~\ref{sec:lmn}) is not formalized. \texttt{RozierTerracol54} is the content of \cite[Corollary~5.4]{RT26}, and \texttt{CensoPropio3e11} is the hypothesis of Theorem~\ref{thm:D}. Theorem~\ref{thm:C} is \lean{region_del_fallo} in \lean{OcelotasCotaGeneral}, and Proposition~\ref{prop:tail} is in \lean{OcelotasColaParametrica}.

\emph{Axioms.} The two external results are hypotheses of the theorems, not axioms. For the proof of Theorem~\ref{thm:A}, the $650$ new modules contain $26\,022$ \texttt{\#print axioms} queries: $26\,013$ report exactly \texttt{propext}, \texttt{Classical.choice} and \texttt{Quot.sound}, five report \texttt{propext} and \texttt{Quot.sound}, and four report only \texttt{propext}. None reports \texttt{sorryAx}, the sources contain no \texttt{native\_decide}, \texttt{implemented\_by}, \texttt{extern} or declared axiom, and the $38$ modules of the common base were checked by hash, with $156$ further queries. For Theorem~\ref{thm:B}, $14$ new modules contain $133$ queries, and for Theorem~\ref{thm:D}, $20$ new modules contain $96$ queries, all with the same three axioms at most. Most queries concern individual pieces of the certificates; they are not independent results.

\emph{Cost.} The proof of Theorem~\ref{thm:A} took about $11.3$ hours of Lean processing in total (several modules compiled in parallel), and its compiled objects occupy $11.1$~GiB; the window $W_2$ accounts for most of both. The audit runs for Theorems~\ref{thm:B} and~\ref{thm:D}, which reused the common base, took about $5$ and $10$ minutes.

\emph{Independent checks.} Separate programs, which do not import the generators of the certificates, recompute the tables and the pieces with rational series for $\operatorname{artanh}$ and exact integer powers, rebuild every leaf of the windows from direct iteration of $T$ on representatives, check that the leaves are disjoint and cover the windows, and recompute the hashes of all sources, compiled objects and audit records. These checks found no discrepancy. They are finite checks of the interpretation; the universal inferences are the theorems proved in Lean.

Lean certifies the displayed statements; whether they express the theorems stated in the introduction is for the reader to check against the definitions above. The computation of Section~\ref{sec:census} is not part of the Lean development.

\section{Remarks}\label{sec:remarks}

\begin{remark}[Orbits without a first crossing]\label{rem:infinite}
If $\tau(n)=\infty$, then $3^{S_j(n)}\ge2^{j}$ for all $j\ge1$, so $T^{j}(n)\ge n$ for all $j$ by Lemma~\ref{lem:affine}. If the orbit of such an $n$ were eventually periodic, then, as in the proof of Corollary~\ref{cor:cycles}, each period would multiply the coefficient by some $3^{K}/2^{L}<1$, and the coefficient would eventually fall below $1$. Hence the orbit of $n$ is unbounded. Our theorems say nothing about such $n$.
\end{remark}

\begin{remark}[Cost of the certificates]\label{rem:cost}
The tail (Proposition~\ref{prop:tail}) and the table (Proposition~\ref{prop:row}) depend on $R$ only through explicit constants, and Proposition~\ref{prop:small} does not depend on $R$. The growth is in the windows. For the analogue of $W_2$ (seeds above $N_0$ with $q\ge190\,537$), the number of nodes was $4\,047$, $67\,649$, $670\,749$, $3\,840\,535$, $12\,368\,501$ and $29\,403\,783$ for $R=19,\dots,24$; the first four used the binary sieve and the last two the refined sieve. The ratio between consecutive values decreases, but each new value of $R$ requires a new and larger computation. As for $m$-cycles, extending $R$ one step at a time does not lead to the full conjecture; that would require a reduction from words with many runs to words with fewer, which we do not have. Such a reduction would in particular exclude all nontrivial cycles.
\end{remark}

\begin{remark}[Earlier stages]
The same method, with the bound of Rhin \cite{Rhi87} in place of \cite{LMN95}, first gave the case of two runs, where everything reduces to $n\le50$ and the tightest case is Example~\ref{ex:39}. Intermediate formalizations for $10$, $13$ and $18$ to $23$ runs are subsumed by Theorem~\ref{thm:A}. A first version of Theorem~\ref{thm:B}, with $77$ runs, used a third window of $1\,420\,851$ nodes; the weights of Section~\ref{sec:weights} replace it. With the same weights, the computation up to $2^{36}$ already gives $40$ runs, and up to $2\cdot10^{11}$, $41$.
\end{remark}

\section*{Use of AI tools}
The proofs, the certificates, the formalization, the numerical checks and the first draft of this text were produced with extensive assistance from large language models used as programming and proof assistants: Anthropic's Claude Opus~5.5 and OpenAI's ChatGPT, between 30 September and 12 October 2026, according to the commit history. The author directed the work and is responsible for the content. Because of this, every step is either formalized in Lean or checked by a separate program, as described above.

\section*{Data availability}
The Lean sources, the certificates in JSON format, their generators, the independent checkers, the outputs of the computation of Section~\ref{sec:census} and the audit records are in the public repository \url{https://github.com/theswordofthemorning/conjeturadecollazpython}, directory \path{depuracion/cst20260930}, at commit \texttt{ccedd6cfd3ff}.

\section*{Statements and declarations}
\emph{Funding.} No funding was received for conducting this study. \emph{Competing interests.} The author has no relevant financial or non-financial interests to disclose.

\begin{thebibliography}{99}

\bibitem{Eve77} C.~J.~Everett, \emph{Iteration of the number-theoretic function $f(2n)=n$, $f(2n+1)=3n+2$}, Adv.\ Math.\ 25 (1977), 42--45.

\bibitem{Her23} C.~Hercher, \emph{There are no Collatz $m$-cycles with $m\le91$}, J.~Integer Seq.\ 26 (2023), Article 23.3.5.

\bibitem{Lag85} J.~C.~Lagarias, \emph{The $3x+1$ problem and its generalizations}, Amer.\ Math.\ Monthly 92 (1985), 3--23.

\bibitem{Lag10} J.~C.~Lagarias (ed.), \emph{The Ultimate Challenge: The $3x+1$ Problem}, Amer.\ Math.\ Soc., Providence, RI, 2010.

\bibitem{LMN95} M.~Laurent, M.~Mignotte and Yu.~Nesterenko, \emph{Formes linéaires en deux logarithmes et déterminants d'interpolation}, J.~Number Theory 55 (1995), 285--321.

\bibitem{Lean4} L.~de Moura and S.~Ullrich, \emph{The Lean~4 theorem prover and programming language}, in: Automated Deduction -- CADE 28, Lecture Notes in Comput.\ Sci.\ 12699, Springer, 2021, 625--635.

\bibitem{Mathlib} The mathlib Community, \emph{The Lean mathematical library}, in: Proceedings of the 9th ACM SIGPLAN International Conference on Certified Programs and Proofs (CPP 2020), ACM, 2020, 367--381.

\bibitem{Rhi87} G.~Rhin, \emph{Approximants de Padé et mesures effectives d'irrationalité}, in: Séminaire de Théorie des Nombres, Paris 1985--86, Progr.\ Math.\ 71, Birkhäuser, Boston, 1987, 155--164.

\bibitem{RT26} O.~Rozier and C.~Terracol, \emph{Paradoxical behavior in Collatz sequences}, Discrete Math.\ 349 (2026), Article 115167. doi:10.1016/j.disc.2026.115167. arXiv:2502.00948.

\bibitem{SdW05} J.~L.~Simons and B.~M.~M.~de Weger, \emph{Theoretical and computational bounds for $m$-cycles of the $3n+1$ problem}, Acta Arith.\ 117 (2005), 51--70.

\bibitem{Ste77} R.~P.~Steiner, \emph{A theorem on the Syracuse problem}, in: Proceedings of the 7th Manitoba Conference on Numerical Mathematics and Computing, Congr.\ Numer.\ 20, Utilitas Math., Winnipeg, 1978, 553--559.

\bibitem{Ter76} R.~Terras, \emph{A stopping time problem on the positive integers}, Acta Arith.\ 30 (1976), 241--252.

\end{thebibliography}

\appendix

\section{The table for 24 runs}\label{app:table}

Row $j$ covers $Q_j\le q<Q_{j+1}$, with $Q_{43}=39\,123\,338\,641$. Its upper pair is $(P_j,Q_j)$, with $P_j\log2-Q_j\log3>0$; the lower pair $(p_j,q'_j)$ of Lemma~\ref{lem:onesided} is determined by $q'_j=Q_{j+1}-Q_j$ and $p_j=(P_jq'_j-1)/Q_j$. The column $s_j$ is the exclusive bound of Proposition~\ref{prop:row}, $B_j$ the least integer with $s_j^{256}\le2^{B_j}$, and $C_{24}(B_j)$ the bound of the refined sieve. The last column says how the row is closed: by Proposition~\ref{prop:small} ($n\le N_0$), by the sieve ($C_{24}(B_j)<Q_j$), or by a window of Table~\ref{tab:windows}. The exact rational lower bounds $c_0$ are part of the certificate.

\begin{center}
\footnotesize
\begin{tabular}{@{}rrrrrrl@{}}
\toprule
$j$ & $Q_j$ & $P_j$ & $s_j$ & $B_j$ & $C_{24}(B_j)$ & closed by \\
\midrule
0 & 1 & 2 & 84 & 1\,637 & 632\,735 & $n\le N_0$ \\
1 & 3 & 5 & 142 & 1\,831 & 735\,348 & $n\le N_0$ \\
2 & 5 & 8 & 461 & 2\,266 & 894\,630 & $n\le N_0$ \\
3 & 17 & 27 & 623 & 2\,377 & 954\,317 & $n\le N_0$ \\
4 & 29 & 46 & 960 & 2\,537 & 1\,002\,885 & $n\le N_0$ \\
5 & 41 & 65 & 2\,094 & 2\,825 & 1\,163\,984 & $n\le N_0$ \\
6 & 94 & 149 & 2\,561 & 2\,899 & 1\,165\,529 & $n\le N_0$ \\
7 & 147 & 233 & 3\,294 & 2\,992 & 1\,212\,574 & $n\le N_0$ \\
8 & 200 & 317 & 4\,617 & 3\,117 & 1\,282\,331 & $n\le N_0$ \\
9 & 253 & 401 & 7\,717 & 3\,306 & 1\,320\,577 & $n\le N_0$ \\
10 & 306 & 485 & 23\,478 & 3\,717 & 1\,511\,292 & $n\le N_0$ \\
11 & 971 & 1\,539 & 24\,526 & 3\,733 & 1\,512\,592 & $n\le N_0$ \\
12 & 1\,636 & 2\,593 & 25\,671 & 3\,750 & 1\,512\,600 & $n\le N_0$ \\
13 & 2\,301 & 3\,647 & 26\,928 & 3\,768 & 1\,520\,950 & $n\le N_0$ \\
14 & 2\,966 & 4\,701 & 28\,315 & 3\,787 & 1\,521\,014 & $n\le N_0$ \\
15 & 3\,631 & 5\,755 & 29\,852 & 3\,806 & 1\,546\,153 & $n\le N_0$ \\
16 & 4\,296 & 6\,809 & 31\,566 & 3\,827 & 1\,546\,163 & $n\le N_0$ \\
17 & 4\,961 & 7\,863 & 33\,489 & 3\,849 & 1\,578\,228 & $n\le N_0$ \\
18 & 5\,626 & 8\,917 & 35\,661 & 3\,872 & 1\,589\,573 & $n\le N_0$ \\
19 & 6\,291 & 9\,971 & 38\,135 & 3\,897 & 1\,594\,697 & $n\le N_0$ \\
20 & 6\,956 & 11\,025 & 40\,977 & 3\,923 & 1\,627\,767 & $n\le N_0$ \\
21 & 7\,621 & 12\,079 & 44\,277 & 3\,952 & 1\,629\,138 & $n\le N_0$ \\
22 & 8\,286 & 13\,133 & 48\,156 & 3\,983 & 1\,629\,138 & $n\le N_0$ \\
23 & 8\,951 & 14\,187 & 52\,778 & 4\,017 & 1\,629\,142 & $n\le N_0$ \\
24 & 9\,616 & 15\,241 & 58\,383 & 4\,054 & 1\,629\,988 & $n\le N_0$ \\
25 & 10\,281 & 16\,295 & 65\,320 & 4\,095 & 1\,642\,957 & $n\le N_0$ \\
26 & 10\,946 & 17\,349 & 74\,126 & 4\,142 & 1\,710\,838 & $n\le N_0$ \\
27 & 11\,611 & 18\,403 & 85\,678 & 4\,195 & 1\,714\,521 & $n\le N_0$ \\
28 & 12\,276 & 19\,457 & 101\,495 & 4\,258 & 1\,728\,196 & $n\le N_0$ \\
29 & 12\,941 & 20\,511 & 124\,475 & 4\,333 & 1\,762\,738 & $n\le N_0$ \\
30 & 13\,606 & 21\,565 & 160\,905 & 4\,428 & 1\,836\,269 & $n\le N_0$ \\
31 & 14\,271 & 22\,619 & 227\,483 & 4\,556 & 1\,847\,365 & $n\le N_0$ \\
32 & 14\,936 & 23\,673 & 388\,044 & 4\,753 & 1\,963\,811 & $n\le N_0$ \\
33 & 15\,601 & 24\,727 & 1\,319\,074 & 5\,205 & 2\,159\,221 & $n\le N_0$ \\
34 & 47\,468 & 75\,235 & 2\,195\,861 & 5\,393 & 2\,243\,510 & $n\le N_0$ \\
35 & 79\,335 & 125\,743 & 6\,548\,931 & 5\,797 & 2\,378\,451 & $W_1$ \\
36 & 190\,537 & 301\,994 & 372\,049\,746 & 7\,289 & 3\,013\,116 & $W_2$ \\
37 & 10\,781\,274 & 17\,087\,915 & 1\,966\,087\,807 & 7\,904 & 3\,263\,200 & sieve \\
38 & 64\,497\,107 & 102\,225\,496 & 2\,747\,757\,915 & 8\,028 & 3\,338\,902 & sieve \\
39 & 118\,212\,940 & 187\,363\,077 & 4\,561\,172\,410 & 8\,215 & 3\,425\,372 & sieve \\
40 & 171\,928\,773 & 272\,500\,658 & 13\,413\,694\,944 & 8\,613 & 3\,562\,837 & sieve \\
41 & 397\,573\,379 & 630\,138\,897 & 226\,749\,896\,056 & 9\,657 & 4\,007\,065 & sieve \\
42 & 6\,586\,818\,670 & 10\,439\,860\,591 & 2\,370\,809\,334\,049 & 10\,524 & 4\,417\,608 & sieve \\
\bottomrule
\end{tabular}
\end{center}

\end{document}
